% !PRINTED 150
\weberpage{150}
The difference
\[
S - P = S'
\]
is therefore again a symmetric function whose highest term is
lower than that of $S$, and which by hypothesis we can represent
rationally by the fundamental functions.
Since $P$ is represented in the same way, the goal is attained. It is
also evident here once again that no fractions are introduced by the
representation by means of the fundamental functions, if originally
none are contained in $S$.

This method of computing symmetric functions at the same
time gives us information about the degree of the integral rational
function of $a_1,\ a_2 \ldots a_n$ by which a given
symmetric function of the $\alpha_1,\ \alpha_2 \ldots \alpha_n$ is expressed.
Namely, if $(\nu_1,\ \nu_2 \ldots \nu_n)$ is the order of the symmetric
function (measured by the sequence of exponents of the first
term), then $\nu_1$ is the degree of $P$ with respect to all of
$a_1, a_2 \ldots a_n$, and in the developed expression for $S$ no
term of higher degree can occur; nor can the term $P$ be
destroyed by any of the following terms, since $P$ is completely
determined by the order of $S$, and the order of $S'$
and of all the following symmetric functions is lower than
the order of $S$.

If, then, in place of $a_1, a_2 \ldots a_n$ we put:
\[
\frac{a_1}{a_0},\quad \frac{a_2}{a_0} \ldots \frac{a_n}{a_0},
\]
then $a_0^{\nu_1}\,S$ is an \gesperrt{integral homogeneous function
of degree} $\nu_1$ \gesperrt{in} $a_0,\ a_1 \ldots a_n$, and one cannot turn $S$ into an integral
function of the $a$ by multiplication with a lower power of $a_0$.

\begin{center}
\S.\,46.

\gesperrt{Discriminants.}
\end{center}

The theorems obtained in the foregoing teach us to express
symmetric functions of the roots of an equation rationally by
the coefficients of the equation; the results are
at first derived under the assumption that the roots,
or, what by the third section amounts to the same thing,
%CONT>>
